\section{Architecture de l'agent}\label{sec:archi}

L'agent est un graphe LangGraph (\texttt{agent/agent\_graph.py}). Chaque nœud lit l'état partagé et ne renvoie que les champs
qu'il modifie. La figure~\ref{fig:graphe} montre les nœuds et les transitions.

\begin{figure}[H]
\centering
\begin{tikzpicture}[
  node distance=7mm and 9mm,
  box/.style={draw, rounded corners=2pt, minimum width=22mm, minimum height=8mm, align=center, font=\small},
  llm/.style={box, fill=blue!8},
  code/.style={box, fill=gray!10},
  todo/.style={box, dashed, fill=white},
  arr/.style={-{Stealth[length=2mm]}, thick},
  lab/.style={font=\scriptsize, fill=white, inner sep=1pt}]
  \node[llm]  (an)  {Analyse};
  \node[code, above right=4mm and 12mm of an] (cl) {Clarifier\\{\scriptsize (au plus 2 fois)}};
  \node[code, below right=4mm and 12mm of an] (re) {Refuser};
  \node[code, right=34mm of an] (rt) {Récupérer\\{\scriptsize (G)}};
  \node[llm, right=of rt] (gr) {Évaluer};
  \node[llm, below=12mm of gr] (rf) {Reformuler};
  \node[code, left=of rf] (rm) {Récupérer+};
  \node[llm, right=of gr] (aw) {Répondre};
  \node[todo, below=12mm of aw] (vf) {Vérifier\\{\scriptsize (hors graphe)}};
  \draw[arr] (an) -- node[lab, above, sloped] {ambiguë} (cl);
  \draw[arr] (an) -- node[lab, below, sloped] {hors champ} (re);
  \draw[arr] (an) -- node[lab, above] {claire} (rt);
  \draw[arr] (rt) -- (gr);
  \draw[arr] (gr) -- node[lab, above] {suffisant} (aw);
  \draw[arr] (gr) -- node[lab, right] {faible} (rf);
  \draw[arr] (rf) -- (rm);
  \draw[arr] (rm) -- node[lab, left] {reprises $<2$} (gr);
  \draw[arr] (rm.south) -- ++(0,-5mm) -| node[lab, pos=0.25, below] {reprises $=2$ (sans réévaluer)} (aw.south east);
  \draw[arr, dashed] (aw) -- (vf);
\end{tikzpicture}
\caption{Graphe de l'agent. En bleu, les nœuds qui appellent le LLM ; en gris, le code déterministe ; en pointillés, le
vérificateur de support, étalonné (section~\ref{sec:generation}) mais pas encore branché dans le graphe.}
\label{fig:graphe}
\end{figure}

\subsection{Les nœuds}
\begin{description}
  \item[Analyse.] Le LLM détecte la langue, classe l'intention (claire, ambiguë, hors champ) et liste les codes qui
        peuvent s'appliquer. Le prompt contient une description d'une ligne de chacun des quatorze codes, rédigée à
        partir de la structure des PDF (tables des matières, titres et chapitres).
  \item[Clarifier / Refuser.] Une question ambiguë reçoit une question de précision ; la réponse de l'usager repasse par
        l'Analyse avec toute la transcription. Une question hors champ est refusée sans recherche.
  \item[Récupérer (G).] La récupération de v1.1 sur la question brute (A2), fusionnée par rang réciproque avec la même
        recherche filtrée sur le code proposé par l'Analyse lorsque celle-ci n'en propose qu'un. Le code est donc un
        indice « doux » : il ne peut pas faire disparaître les résultats non filtrés.
  \item[Évaluer.] Le LLM lit les six articles (500 premiers caractères de chacun) et juge si l'un d'eux énonce la règle
        demandée ; sinon il nomme, en français juridique, la règle manquante.
  \item[Reformuler et Récupérer+.] Le LLM écrit une à trois nouvelles requêtes en français, différentes des précédentes et
        réutilisant les termes vus dans les articles ; les nouveaux résultats sont fusionnés avec les anciens.
  \item[Répondre.] Le LLM écrit de une à six affirmations ; chacune est une phrase qui cite les identifiants des articles
        qui l'énoncent. Un contrôle en code vérifie que chaque identifiant cité figure bien dans les preuves.
\end{description}

\subsection{L'état}
L'état transporte la question et le texte complet de la conversation (\texttt{full\_question}, utilisé par tous les nœuds
après une clarification), la langue, l'intention, les codes, les six documents et les preuves (texte des articles indexé
par identifiant), le nombre de reprises, les requêtes déjà essayées, le texte de la règle manquante, les affirmations
structurées, les identifiants cités, les problèmes détectés et une liste \texttt{degraded} qui enregistre les replis
silencieux (par exemple une Analyse en échec).

\subsection{Boucles bornées}
Deux boucles existent et toutes deux ont une borne dans l'état, non dans la bibliothèque : au plus deux reprises de
recherche, et au plus deux questions de précision. Après la dernière reformulation, le graphe répond sans réévaluer :
une évaluation de plus ne pourrait plus rien changer et coûterait un appel LLM.
